% Generated by scripts/tikz_reviews.py from the figure manifests.
\ReviewFigure{PCA directions and projected variance}
{figures/ml/variance.png}
{figures/tikz/ml-pca-variance.pdf}
{Principal directions are eigenvectors of the empirical covariance; the chapter writes its eigenvalues as sigma squared.}
{This centered illustrative cloud has exact reflection symmetries in its principal coordinates. The solid ellipse has semiaxes equal to the empirical standard deviations, using the chapter normalization 1/n. Variance is sigma squared, whereas the indicated geometric lengths are sigma. The point locations are illustrative, not recovered measurements.}
{ml-pca-variance}
\ReviewFigure{Nearest-centroid assignment and Voronoi cells}
{figures/ml/voronoi.png}
{figures/tikz/ml-voronoi.pdf}
{For fixed centroids, the assignment step of k-means selects the nearest centroid in Euclidean distance.}
{Cell edges are actual perpendicular bisectors of the displayed centroids. Their meeting point is equidistant from all three. Colored sample points lie in the corresponding nearest-centroid cells. The diagram illustrates an assignment step and does not assert that these centroids are already cluster means.}
{ml-voronoi}
\ReviewFigure{A joint probability model for input and output}
{figures/ml/proba-model.pdf}
{figures/tikz/ml-joint-model.pdf}
{Supervised observations are pairs (x\_i,y\_i) drawn from a joint law pi on the product of input and output spaces.}
{Contours represent an illustrative bivariate Gaussian density, not an asserted fit to the original hand-drawn points. The dashed projections identify the two coordinates of one observation. The general discussion in the chapter is not restricted to Gaussian models.}
{ml-joint-model}
\ReviewFigure{Conditional expectation as the mean of a slice}
{figures/ml/cond-expect.pdf}
{figures/tikz/ml-conditional-expectation.pdf}
{Under squared loss, the optimal predictor is the conditional expectation of the output given the input.}
{An illustrative model Y=m(X)+epsilon with centered Gaussian noise is used. The left curves show constant conditional-density levels around m(x); the right density is normalized at x=x0. Its mean is m(x0), which need not be linear in x. No distributional assumption is added to the theorem.}
{ml-conditional-expectation}
\ReviewFigure{An affine fit and a nonlinear regression function}
{figures/ml/linear-fit.png}
{figures/tikz/ml-linear-fit.pdf}
{Linear regression restricts the predictor to a finite-dimensional model, even when the conditional expectation is nonlinear.}
{The red line is the least-squares affine fit to the actual displayed illustrative points. The dashed curve is the model conditional mean used to construct those points. Residuals are vertical output errors, not orthogonal distances to the line.}
{ml-linear-fit}
\ReviewFigure{Four nearest neighbors of a query point}
{figures/ml/clustering.png}
{figures/tikz/ml-nearest-neighbors.pdf}
{The R-nearest-neighbor rule sorts training points by their Euclidean distance from the query.}
{The arrows identify the four nearest observations in increasing distance order. The dashed circle has radius equal to the fourth distance. All remaining points are farther away. This is a nearest-neighbor query, not the centroid assignment used in k-means.}
{ml-nearest-neighbors}
\ReviewFigure{Logistic probability and transition width}
{figures/ml/logistic-1d.png}
{figures/tikz/ml-logistic-one-dimension.pdf}
{The probability of class +1 is theta(beta x), with theta(s)=1/(1+exp(-s)).}
{The increasing sign convention follows the current equation. Two positive slopes show the same probability-one-half boundary at x=0; the larger slope gives a narrower transition. The vertical axis is a probability in [0,1], not a class label in \{-1,+1\}.}
{ml-logistic-one-dimension}
\ReviewFigure{The same decision boundary at two parameter norms}
{figures/ml/logistic-2d.png}
{figures/tikz/ml-logistic-two-dimensions.pdf}
{The parameter direction determines a hyperplane normal; its norm controls the width of the logistic probability transition.}
{Both panels use the same unit normal v and beta=t v. The shaded bands delimit probabilities from 0.1 to 0.9, so their normal width is 2 log(9)/t. Changing t leaves the decision boundary unchanged; it does not change whether a data set is linearly separable.}
{ml-logistic-two-dimensions}
\ReviewFigure{Rows are observations; columns are features}
{figures/ml/pca-maths/pca-maths-1.pdf}
{figures/tikz/ml-pca-data-matrix.pdf}
{The PCA proof centers the data and writes the n-by-p matrix X with observations in its rows.}
{The reconstruction follows the current n-observations, p-features convention. An observation x\_i is a column vector in R\textasciicircum{}p; the corresponding matrix row is its transpose. Centering is explicit, and the covariance is X transpose X divided by n.}
{ml-pca-data-matrix}
\ReviewFigure{Linear compression followed by reconstruction}
{figures/ml/pca-maths/pca-maths-2.pdf}
{figures/tikz/ml-pca-compression-reconstruction.pdf}
{The variational PCA problem compares rank-at-most-k linear reconstructions R S transpose.}
{Both R and S have p rows and k columns. Compression applies S transpose; reconstruction applies R, so the product is a p-by-p map of rank at most k. Orthogonality is not assumed for arbitrary R and S; it is justified later when reducing to orthogonal projections.}
{ml-pca-compression-reconstruction}
\ReviewFigure{The polytope bounding the PCA objective}
{figures/ml/pca-maths/pca-maths-3.pdf}
{figures/tikz/ml-pca-linear-program.pdf}
{The proof bounds the trace by maximizing a decreasingly weighted sum over 0 <= beta <= 1 and sum beta <= k.}
{The left feasible set is the triangle for p=2, k=1. The right is the unit cube with the corner (1,1,1) cut off for p=3, k=2. The marked maximizers are (1,0) and (1,1,0), respectively, for ordered nonnegative eigenvalues. The three-dimensional shaded face is sum beta=2.}
{ml-pca-linear-program}
\ReviewFigure{Row norms of the rotated basis}
{figures/ml/pca-maths/pca-maths-4.pdf}
{figures/tikz/ml-pca-row-norms.pdf}
{The PCA proof defines B=V transpose S and beta\_i as the squared norm of its ith row.}
{The matrix B has p rows and k orthonormal columns. The vector b\_i is in R\textasciicircum{}k and appears transposed as a row. Its squared row norms sum to k, the squared Frobenius norm of B. These dimensions follow the current proof rather than ambiguous letters in the scan.}
{ml-pca-row-norms}
\ReviewFigure{Completing the columns to an orthogonal matrix}
{figures/ml/pca-maths/pca-maths-5.pdf}
{figures/tikz/ml-pca-orthogonal-extension.pdf}
{An orthogonal completion proves that every squared row norm beta\_i is at most one.}
{The first k columns are B; the remaining p-k columns form B perpendicular. Every full row of the square orthogonal matrix has norm one. Restricting a row to the first k entries cannot increase its norm. When k=p, the complementary block is empty.}
{ml-pca-orthogonal-extension}
